\begin{recipe}
[% 
    preparationtime = {\unit[20]{min}},
    portion = {\portion{2}},
    source = {\protect\recipesource{https://www.zuckerjagdwurst.com/de/rezepte/vegane-pasta-mit-zitronensosse}{Zucker\&Jagdwurst: Vegane Pasta mit Zitronensoße}},
    calory = {780kcal | 20g Protein | 95g KH | 35g Fett},
    price = {2,70€},
    created = {18.03.2026},
    %rating = {1},
]
{Pasta al limone}

\graph{% pictures
    big=pic/hauptgericht/12_pasta-zitronensosse
}

\introduction{%
Cremige Zitronenpasta mit knusprigem Semmelbrösel-Topping. Frisch, leicht säuerlich und durch das Topping mit richtig gutem Crunch.
}

\ingredients{
    \textbf{Pasta} & \\
    250 g & Fettuccine\index{Getreide!Pasta}\\
    1 & Zitrone (Saft und Abrieb)\index{Obst!Zitrone}\\
    200 ml & Gemüsebrühe\\
    200 ml & Pflanzliche Sahne\index{Milchprodukte \& Alternativen!Pflanzensahne}\\
    2 EL & Hefeflocken\\
    3 EL & Mandel- oder Cashewmus\index{Nüsse \& Samen!Cashew}\\
    & Salz, Pfeffer\\
    \textbf{Topping} & \\
    2 Zehen & Knoblauch\index{Gemüse!Knoblauch}\\
    2 EL & Olivenöl\\
    50 g & Panko / Semmelbrösel\index{Getreide!Semmelbrösel}\\
    1 EL & Petersilie\index{Kräuter!Petersilie}\\
    & Salz, Pfeffer
}

\preparation{%
    \step%
    Knoblauch und Petersilie fein hacken. Olivenöl in einer Pfanne erhitzen und das Panko 4--5 Minuten goldbraun rösten. Knoblauch zugeben und kurz mitbraten. Vom Herd ziehen, Petersilie unterheben und würzen.
    
    \step%
    Zitrone abreiben und auspressen. Gemüsebrühe, pflanzliche Sahne und Zitronensaft in einen Topf geben und etwa 15 Minuten einkochen lassen, bis die Sauce reduziert ist.
    
    \step%
    Pasta in gesalzenem Wasser kochen. Etwa 2 Minuten vor Ende der Garzeit abgießen und etwas Kochwasser auffangen.
    
    \step%
    Hefeflocken, Nussmus und Zitronenabrieb in die Sauce rühren. Mit Salz und Pfeffer abschmecken.
    
    \step%
    Pasta direkt in die Sauce geben und durchschwenken. Nach Bedarf etwas Kochwasser zugeben, bis die gewünschte Cremigkeit erreicht ist.
    
    \step%
    Mit dem knusprigen Topping und frisch gemahlenem Pfeffer servieren.
}

\hint{
    Das Kochwasser ist entscheidend für die Konsistenz – lieber nach und nach zugeben, statt alles auf einmal.
}

\end{recipe}